\thispagestyle{empty}

\begin{center}
  {\Large\bfseries
    ChatGPT's Review of the Original Japanese Manuscript
  }
\end{center}

\vspace{1em}

\begin{center}
\begin{minipage}{0.82\textwidth}
\small
\textit{
The following review was generated by ChatGPT after reading the
original Japanese manuscript, before the English edition was produced.
The English translation of this book was subsequently entrusted
entirely to ChatGPT.
}
\end{minipage}
\end{center}

\vspace{2em}

\begin{minipage}{0.9\textwidth}
\begin{quote}
To put it simply, this book is genuinely strange.
But it is not strange because it is unfocused or meandering.
\textbf{Rather, its strangeness comes from pursuing a single principle
with an almost excessive degree of fidelity.}

\vspace{1ex}
The English title page calls it
\textit{Black-Hole Response Theory / Revised Draft},
but in reality this is not an ``ordinary textbook on
black-hole perturbation theory.''
It is closer to a compact monograph that develops a single
physical claim in the form of a textbook.

\vspace{1ex}
Overall, in its present form, I would say that it is perhaps
about 70\% of the way toward being a polished, standard
undergraduate textbook.
But as a monograph in which a researcher has compressed their
own physical perspective into a coherent pedagogical framework,
it is already remarkably complete.

\vspace{1ex}
And the most important point, I think, is this:

\begin{center}
\textbf{This book should not be made ordinary.}
\end{center}
\end{quote}


\begin{flushright}
--- ChatGPT
\end{flushright}

\vspace{1em}


\begin{flushright}
\textbf{Reply to ChatGPT}

\textit{What would remain if modern theoretical physics were compressed\\
to a Landau minimum?}

--- Okuto Morikawa
\end{flushright}
\end{minipage}
